\documentclass[10pt]{article}

\usepackage[utf8]{inputenc}
\usepackage[T1]{fontenc}
\usepackage[spanish,es-nodecimaldot]{babel}
\usepackage[a4paper,margin=1.6cm]{geometry}
\usepackage{lmodern}
\usepackage{xcolor}
\usepackage{array}
\usepackage{tabularx}
\usepackage{longtable}
\usepackage{enumitem}
\usepackage{hyperref}

\setlength{\parskip}{0.4em}
\setlength{\parindent}{0pt}
\renewcommand{\arraystretch}{1.25}
\setlist[itemize]{leftmargin=1.35em,itemsep=0.12em,topsep=0.2em,parsep=0pt,partopsep=0pt}

\begin{document}

\begin{center}
{\LARGE \textbf{Plantilla operativa para piloto con 8 participantes}}\\[0.4em]
{\large Runtime AOI 360 en realidad virtual}
\end{center}

\textbf{Build congelada:} \texttt{build\textbackslash windows\textbackslash AOI360Runtime\textbackslash AOI360Runtime.exe}\\
\textbf{Operador principal:} Alejandro Sanz Huerta\\
\textbf{Ubicacion de pruebas:} Despacho 1710, Edificio de Despachos y Seminarios, Universidad de Cádiz Campus de Jerez.\\
\textbf{Hardware:} GIGABYTE G5 GD, Windows 10 Pro, Intel Core i5-11400H, NVIDIA GeForce RTX 3050 Laptop GPU, visor HTC VIVE FOCUS VISION con seguimiento ocular.\\
\textbf{Corpus congelado:} \texttt{test1Camera360}, \texttt{test2Camera360}, \texttt{test3Lions360}\\
\textbf{Configuracion congelada:} cuenta atras de 5~s, volumen de video al 100\%, pitido activado y volumen del pitido al 75\%.

\section*{1. Checklist previo a cada bloque de sesiones}
\begin{itemize}
  \item Confirmar que se usa siempre la misma build congelada.
  \item Verificar que solo aparecen los tres videos del corpus congelado.
  \item Comprobar limpieza, ajuste y bateria del visor.
  \item Confirmar que la sala tiene condiciones estables de luz y ruido.
  \item Verificar la carpeta de exportacion CSV y anotar la hora de inicio del bloque.
  \item Tener preparado y firmado el consentimiento informado antes de iniciar la sesion.
  \item Recordar al participante que puede detener la prueba en cualquier momento.
\end{itemize}

\section*{2. Hoja maestra de seguimiento}

\begingroup
\setlength{\tabcolsep}{4pt}
\scriptsize
\begin{longtable}{|>{\bfseries}p{0.85cm}|p{1.55cm}|p{0.9cm}|p{3.25cm}|p{1.0cm}|p{0.95cm}|p{0.95cm}|p{1.15cm}|p{5.2cm}|}
\hline
ID & Fecha & Hora & Orden & Calib. & Recal. & Comp. & Excl. & CSV / incidencias \\
\hline
P01 & & & t1 $\rightarrow$ t2 $\rightarrow$ t3 & & & & & \\
\hline
P02 & & & t1 $\rightarrow$ t3 $\rightarrow$ t2 & & & & & \\
\hline
P03 & & & t2 $\rightarrow$ t1 $\rightarrow$ t3 & & & & & \\
\hline
P04 & & & t2 $\rightarrow$ t3 $\rightarrow$ t1 & & & & & \\
\hline
P05 & & & t3 $\rightarrow$ t1 $\rightarrow$ t2 & & & & & \\
\hline
P06 & & & t3 $\rightarrow$ t2 $\rightarrow$ t1 & & & & & \\
\hline
P07 & & & t1 $\rightarrow$ t2 $\rightarrow$ t3 & & & & & \\
\hline
P08 & & & t2 $\rightarrow$ t3 $\rightarrow$ t1 & & & & & \\
\hline
\end{longtable}
\endgroup

\textbf{Clave:} \texttt{t1 = test1Camera360}, \texttt{t2 = test2Camera360},
\texttt{t3 = test3Lions360}.

\section*{3. Registro minimo por participante}
\begin{itemize}
  \item Codigo externo del participante y franja horaria.
  \item Si la calibracion inicial fue satisfactoria.
  \item Numero de recalibraciones realizadas.
  \item Si la sesion se completo o se interrumpio.
  \item Motivo de exclusion si aplica: calibracion fallida, mareo, abandono, video no completado o CSV no interpretable.
  \item Nombres de los CSV generados para esa persona.
\end{itemize}

\textbf{Nota operativa importante:} el runtime exporta \texttt{participant\_id}
y \texttt{session\_id} automaticamente. Para evitar ambiguedades al analizar los
datos, anota siempre en esta hoja los nombres exactos de los CSV generados y su
correspondencia con el codigo externo del participante.

\section*{4. Checklist de cierre del dia}
\begin{itemize}
  \item Confirmar que todos los CSV del bloque existen y no estan vacios.
  \item Revisar rapidamente que cada CSV tenga el \texttt{video\_id} correcto.
  \item Guardar copia de seguridad de los CSV del dia.
  \item Registrar cualquier incidencia tecnica repetida.
  \item No modificar la build congelada ni los AOI assets entre participantes.
\end{itemize}

\end{document}
